\subsection{Detailed Experimental Design}
\label{app:detailed-experimental-design}

Unless stated otherwise, we use the default hyperparameters from the original
environment or algorithm implementations used in the corresponding papers.
This appendix only records the choices that are specific to our convergence
experiments.

\subsubsection{Office-World}
\label{app:design-office}

Office-World uses tabular value functions.
The primitive environment is \texttt{Office-v0}; the evaluation tasks are Coffee, Patrol, Coffee and Mail, and Long.
The primitive library is
\[
    \mathcal{B}
    =
    \{0,1\}
    \cup
    \{p_\varphi:\varphi\in\{A,B,C,D,c,d,m,o,tm,to\}\}
    \cup
    \{c_d\},
\]
so each run trains $13$ tabular world value functions.
The propositions $A,B,C,D$ denote rooms, $c$ coffee, $m$ mail, $o$ office, $tm$ and $to$ the dynamic mail/person availability propositions, and $d$ the decoration constraint.
We use Q-learning with learning rate $0.1$, discount $0.9$, $\epsilon$-greedy exploration with $\epsilon=0.5$, and zero initialization.
We run $3$ independent seeds.
Evaluation uses $50$ episodes per task, method, and checkpoint, with horizon $1000$ and evaluation discount $0.9$.
We compare Original, Univ./empty, and Base tasks; the first two use only the learned $0$ and $1$ primitives, while Base tasks can use all primitives in $\mathcal{B}$.

\subsubsection{Safety-Gym}
\label{app:design-safety-gym}

Safety-Gym is the continuous-control function-approximation experiment.
The primitive environment is \texttt{Safety-v0}; evaluation uses \texttt{Safety-Task-1-v0} through \texttt{Safety-Task-6-v0}.
The learned primitive library is
\[
    \{0,1,p_{\textsc{buttons}},p_{\textsc{goal}},
    p_{\textsc{hazards}},c_{\textsc{hazards}}\},
\]
so each run trains $6$ world value functions.
Each primitive is trained with TD3.
The actor and critic use the default TD3 multi-input policy with hidden layers $[2024,2024,2024]$; Gaussian action noise has standard deviation $0.2$.
The main training parameters are learning rate $10^{-6}$, discount $0.99$, batch size $32$, replay buffer size $10^6$, learning starts after $1000$ transitions, training frequency $50$, and $50$ gradient steps per update.
Checkpoints are evaluated at $10^5,2\cdot10^5,4\cdot10^5,6\cdot10^5,8\cdot10^5,10^6$ environment interactions per primitive.
We run $3$ independent seeds.
Evaluation uses $100$ episodes per task, method, and checkpoint, with horizon $1000$ and evaluation discount $0.99$.
We compare Original, Univ./empty, and Base tasks.

\subsubsection{Four Rooms}
\label{app:design-four-rooms}

Four Rooms uses tabular goal-oriented Q-learning.
In the convergence experiment we use the four-room configuration with $4$ terminal goals.
The trained value functions are the universal task, the empty task, and the $\lceil\log_2 4\rceil=2$ binary basis tasks.
Training uses sparse rewards with shared terminal states.
Goal-oriented Q-learning uses the implementation defaults: discount $1$, learning rate $1$, $\epsilon$-greedy exploration with $\epsilon=1$, and a maximum of $100$ steps per episode.
The training budgets range between $10$ and $5000$ episodes per trained task, as we empirically see that this range is enough for convergence.
Evaluation samples tasks proportionally by task cardinality: with $4$ goals this gives $15$ sampled tasks, of which the empty and universal tasks are excluded from the composed-task evaluation, leaving $13$ tasks.
For each budget and composed task, we evaluate $1000$ rollouts for Univ./empty and $1000$ rollouts for Boolean composition, with horizon $100$.

\subsubsection{Boxman}
\label{app:design-boxman-sts}

Boxman STS is the image-based function-approximation experiment.
We train four DQN value functions per independent run: \texttt{on}, \texttt{off}, \texttt{blue}, and \texttt{square}.
These are evaluated on five composed tasks: $B$, $S$, $B+S$, $B.S$, and $B\operatorname{xor}S$.
We run $3$ independent full runs.
Each DQN receives the RGB observation concatenated with the RGB goal image, giving $6$ input channels.
The network is the original Boxman DQN architecture: convolutional layers $32$ filters of size $8$ and stride $4$, $64$ filters of size $4$ and stride $2$, $64$ filters of size $3$ and stride $1$, followed by a fully connected layer of width $512$ and a linear action head.
We use Adam with learning rate $10^{-4}$, discount $0.99$, batch size $128$, replay buffer size $3\cdot10^5$, learning starts after $10^4$ transitions, training frequency $4$, target network update frequency $1000$, and Huber loss with gradient clipping to $[-1,1]$.
Exploration is linearly annealed from $1.0$ to $0.01$ over $10^5$ steps in the convergence runs.
Checkpoints are evaluated at $10^4,2\cdot10^4,\ldots,10^5$ and then every $2\cdot10^4$ steps up to $3\cdot10^5$.
Evaluation uses $100$ episodes per task, method, checkpoint, and completed run, with maximum trajectory length $20$.
We compare the Univ./empty DQN composition against Boolean composition over the blue and square DQNs.

